\Askhsh{23197}{2}
\begin{ylh}{1}
    \item 1.3 Μονότονες συναρτήσεις - Αντίστροφη συνάρτηση
    \item 2.6 Συνέπειες του Θεωρήματος Μέσης Τιμής
    \item 2.7 Τοπικά ακρότατα συνάρτησης.
\end{ylh}
Δίνεται η συνάρτηση $f(x)=x^2-2x$, $x \in \mathbb{R}$.
\begin{alist}
\item Να βρείτε δύο διαφορετικούς αριθμούς $\alpha$, $\beta$ ώστε $f(\alpha)=f(\beta)$. Κατόπιν να αιτιολογήσετε γιατί η συνάρτηση $f$ δεν αντιστρέφεται. \monades{9}
\item Να μελετήσετε τη συνάρτηση, με τη βοήθεια της παραγώγου ή με οποιονδήποτε άλλο τρόπο, ως προς τη μονοτονία και τα ακρότατα. \monades{8}
\item Να σχεδιάσετε τη γραφική παράσταση $C_f$ της $f$. \monades{8}
\end{alist}